% !TEX root = ../tecnologia-en-marcha.tex
% !TEX encoding = UTF-8 Unicode
% =============================================================================
% Título, autores, resumen y palabras clave.
% Texto: primer borrador del artículo (articulo_tecnologia_en_marcha.md), con los
% acentos restaurados.
% =============================================================================

% El título agrega lo que identifica a la investigación (PIA02): la ENIA y las
% personas adultas mayores. La traducción al inglés es una propuesta: la
% plantilla pide evitar traductores automáticos.
\newcommand{\tituloES}{Marco operativo de cumplimiento de la Estrategia Nacional de IA de Costa Rica para el monitoreo domiciliario de personas adultas mayores con sensores PIR}
\newcommand{\tituloEN}{Operational compliance framework for Costa Rica's National AI Strategy in home monitoring of older adults with PIR sensors}

% Una persona autora por bloque: nombre en negrita; afiliación y país; correo;
% ORCID. La plantilla deja lugar para tres.
\newcommand{\autor}[4]{%
  \textbf{#1}\\
  #2\\
  #3\\
  ORCID: #4\par\medskip}

\newcommand{\bloqueAutores}{%
  \autor{Miguel Ángel Méndez Arias}
        {Universidad CENFOTEC. Costa Rica}
        {\pendiente{correo}}
        {\pendiente{ORCID}}
  \autor{Laura María Segura Miranda}
        {Universidad CENFOTEC. Costa Rica}
        {\pendiente{correo}}
        {\pendiente{ORCID}}
  \pendiente{orden de las personas autoras y si hay una tercera (p. ej., quien dirigió la tesis)}\par}

\newcommand{\resumenES}{%
La Estrategia Nacional de Inteligencia Artificial de Costa Rica establece
principios de transparencia, equidad y responsabilidad, pero no especifica
mecanismos técnicos para verificar su cumplimiento en sistemas concretos. Este
artículo presenta y evalúa un marco operativo que traduce esos principios en
una batería de pruebas computacionales y artefactos auditables, aplicada a
sistemas de clasificación de actividades basados en sensores de infrarrojo
pasivo y otros sensores ambientales para el monitoreo domiciliario de personas
adultas mayores. La metodología consistió en traducir tres
principios de la estrategia nacional en nueve requerimientos operativos,
articularlos con el AI Act de la Unión Europea y el NIST AI Risk Management
Framework, e implementarlos como un procedimiento reproducible de seis pasos. El
marco integra explicabilidad local y global mediante SHAP, evaluación de equidad
por subgrupos con umbrales declarados previamente, registro estructurado de
inferencias y generación de artefactos auditables. La aplicación piloto sobre
dos configuraciones del conjunto CASAS produjo expedientes completos: Aruba, con
11\,175 inferencias, y una serie multi-hogar, con 7561 inferencias. En ambos
casos se cubrieron los nueve requerimientos y se obtuvo un veredicto de no
cumplimiento del sistema evaluado. Los resultados muestran que el marco permite
transformar principios normativos en evidencia técnica verificable y
reproducible.}

\newcommand{\palabrasclaveES}{%
Inteligencia artificial; sensores; aprendizaje automático; explicabilidad;
equidad; trazabilidad; auditoría
\revisar{la plantilla pide términos del tesauro de la UNESCO}}

\newcommand{\resumenEN}{%
Costa Rica's National Artificial Intelligence Strategy establishes principles
of transparency, fairness, and accountability, but does not specify technical
mechanisms to verify compliance with them in concrete systems. This article
presents and evaluates an operational framework that translates these
principles into a battery of computational tests and auditable artifacts,
applied to activity classification systems based on passive infrared and other
environmental sensors for home monitoring of older adults. The
methodology translated three principles of the national strategy into nine
operational requirements, aligned them with the European Union AI Act and the
NIST AI Risk Management Framework, and implemented them as a reproducible
six-step procedure. The framework integrates local and global explainability
through SHAP, subgroup fairness evaluation with thresholds declared in advance,
structured inference logging, and generation of auditable artifacts. The pilot
application on two CASAS dataset configurations produced complete evidence
files: Aruba, with 11,175 inferences, and a multi-home series, with 7,561
inferences. In both cases, all nine requirements were covered and the evaluated
system received a non-compliance verdict. The results show that the framework
can transform normative principles into verifiable and reproducible technical
evidence.}

\newcommand{\palabrasclaveEN}{%
Artificial intelligence; sensors; machine learning; explainability; fairness;
traceability; audit}
